% GENERATED FILE. Edit canonical lessons, then run npm run generate:books.
% canonical-source-hash: fnv1a64:810b9214df165834
% canonical-lessons: RU-C49-table, RU-C49-chair, RU-C49-bed, RU-C49-lamp, RU-C49-key

\chapter{Table, Chair, Bed, Lamp, Key}
\label{ch:ru-table-chair-bed-lamp-key}
\hlchaptermodality{49}

\begin{chapteropening}
\textbf{By the end of this chapter:} \emph{I can name a table, a chair, a bed, a lamp and a key.}

Complete the last lesson of chapter 49: I can name a table, a chair, a bed, a lamp and a key.
\end{chapteropening}

\section[\texorpdfstring{\ru{стол} (stol)}{стол (stol)}]{\ru{стол} (stol) — a table}

\label{lesson:RU-C49-table}



\begin{quote}
Before the new one: say the Russian for a window, then the Russian for a door.
\end{quote}

\subsection*{You'll want to know: \ru{стол}}
\textbf{\ru{стол}} (\emph{stol}) — \textquotedblleft{}a table\textquotedblright{}.

\begin{cousinweb}
\textbf{\ru{стол}} $\leftarrow$ Proto-Slavic \emph{*stolŭ} — inherited; Old Russian \emph{stolŭ} also meant \textquotedblleft{}throne, seat\textquotedblright{} (Vasmer). Vasmer compares Gothic \emph{stōls}, \textquotedblleft{}chair\textquotedblright{}, and Lithuanian \emph{stãlas}, \textquotedblleft{}table\textquotedblright{}, but also weighs the root of \emph{\ru{стелить}}, \textquotedblleft{}to spread\textquotedblright{}, so the deeper origin is open.
\end{cousinweb}

\begin{grammarlens}[title={What you've built}]
One more word for this chapter.
\end{grammarlens}

\subsection*{Your turn}
\begin{itemize}
\raggedright
  \item \emph{Say it:} \emph{stol}
  \item \emph{Say it:} \emph{stol}, once more
  \item \emph{Say it:} say \emph{stol}
  \item \emph{Recall:} say the Russian for a window, then the Russian for a door, then say \emph{stol} again
\end{itemize}

\begin{tcolorbox}[breakable,colback=teal!4,colframe=teal!35!black,title={Before you move on}]
What does \ru{стол} mean? (\textquotedblleft{}A table\textquotedblright{}.) Say it once more.
\end{tcolorbox}
\section[\texorpdfstring{\ru{стул} (stul)}{стул (stul)}]{\ru{стул} (stul) — a chair}

\label{lesson:RU-C49-chair}



\begin{quote}
Before the new one: say the Russian for a door, then the Russian for a table.
\end{quote}

\subsection*{You'll want to know: \ru{стул}}
\textbf{\ru{стул}} (\emph{stul}) — \textquotedblleft{}a chair\textquotedblright{}.

\begin{cousinweb}
\textbf{\ru{стул}} — first recorded in 1578, perhaps borrowed from Old Norse \emph{stóll} or Low German \emph{stuhl}, \textquotedblleft{}chair\textquotedblright{} (Vasmer). A High German loan would have begun \emph{\ru{шт}-}.
\end{cousinweb}

\begin{grammarlens}[title={What you've built}]
One more word for this chapter.
\end{grammarlens}

\subsection*{Your turn}
\begin{itemize}
\raggedright
  \item \emph{Say it:} \emph{stul}
  \item \emph{Say it:} \emph{stul}, once more
  \item \emph{Say it:} say \emph{stul}
  \item \emph{Recall:} say the Russian for a door, then the Russian for a table, then say \emph{stul} again
\end{itemize}

\begin{tcolorbox}[breakable,colback=teal!4,colframe=teal!35!black,title={Before you move on}]
What does \ru{стул} mean? (\textquotedblleft{}A chair\textquotedblright{}.) Say it once more.
\end{tcolorbox}
\section[\texorpdfstring{\ru{кровать} (krovát)}{кровать (krovát)}]{\ru{кровать} (krovát) — a bed}

\label{lesson:RU-C49-bed}



\begin{quote}
Before the new one: say the Russian for a table, then the Russian for a chair.
\end{quote}

\subsection*{You'll want to know: \ru{кровать}}
\textbf{\ru{кровать}} (\emph{krovát}) — \textquotedblleft{}a bed\textquotedblright{}.

\begin{cousinweb}
\textbf{\ru{кровать}} $\leftarrow$ Byzantine Greek \emph{krabbátion}, from Greek \emph{krábbatos}, \textquotedblleft{}a pallet, a camp bed\textquotedblright{} (Vasmer).
\end{cousinweb}

\begin{grammarlens}[title={What you've built}]
One more word for this chapter.
\end{grammarlens}

\subsection*{Your turn}
\begin{itemize}
\raggedright
  \item \emph{Say it:} \emph{krovát}
  \item \emph{Say it:} \emph{krovát}, once more
  \item \emph{Say it:} say \emph{krovát}
  \item \emph{Recall:} say the Russian for a table, then the Russian for a chair, then say \emph{krovát} again
\end{itemize}

\begin{tcolorbox}[breakable,colback=teal!4,colframe=teal!35!black,title={Before you move on}]
What does \ru{кровать} mean? (\textquotedblleft{}A bed\textquotedblright{}.) Say it once more.
\end{tcolorbox}
\section[\texorpdfstring{\ru{лампа} (lámpa)}{лампа (lámpa)}]{\ru{лампа} (lámpa) — a lamp}

\label{lesson:RU-C49-lamp}



\begin{quote}
Before the new one: say the Russian for a chair, then the Russian for a bed.
\end{quote}

\subsection*{You'll want to know: \ru{лампа}}
\textbf{\ru{лампа}} (\emph{lámpa}) — \textquotedblleft{}a lamp\textquotedblright{}.

\begin{cousinweb}
\textbf{\ru{лампа}} came from German \emph{Lampe} or French \emph{lampe}, from Latin \emph{lampada} and Greek \emph{lampás}, \textquotedblleft{}a torch\textquotedblright{}; the old form \emph{\ru{лямпа}} came through Polish (Vasmer).
\end{cousinweb}

\begin{grammarlens}[title={What you've built}]
One more word for this chapter.
\end{grammarlens}

\subsection*{Your turn}
\begin{itemize}
\raggedright
  \item \emph{Say it:} \emph{lámpa}
  \item \emph{Say it:} \emph{lámpa}, once more
  \item \emph{Say it:} say \emph{lámpa}
  \item \emph{Recall:} say the Russian for a chair, then the Russian for a bed, then say \emph{lámpa} again
\end{itemize}

\begin{tcolorbox}[breakable,colback=teal!4,colframe=teal!35!black,title={Before you move on}]
What does \ru{лампа} mean? (\textquotedblleft{}A lamp\textquotedblright{}.) Say it once more.
\end{tcolorbox}
\section[\texorpdfstring{\ru{ключ} (klyuch)}{ключ (klyuch)}]{\ru{ключ} (klyuch) — a key}

\label{lesson:RU-C49-key}



\begin{quote}
Before the new one: say the Russian for a bed, then the Russian for a lamp.
\end{quote}

\subsection*{You'll want to know: \ru{ключ}}
\textbf{\ru{ключ}} (\emph{klyuch}) — \textquotedblleft{}a key\textquotedblright{}.

\begin{cousinweb}
\textbf{\ru{ключ}} $\leftarrow$ Proto-Slavic \emph{*ključĭ} — inherited. Vasmer relates it to Greek \emph{kleís}, \textquotedblleft{}key\textquotedblright{}, and Latin \emph{clāvis}, \textquotedblleft{}key\textquotedblright{}, and \emph{claudō}, \textquotedblleft{}I shut\textquotedblright{}. The \emph{\ru{ключ}} that means \textquotedblleft{}a spring\textquotedblright{} may be the same word or another.
\end{cousinweb}

\begin{grammarlens}[title={What you've built}]
One more word for this chapter.
\end{grammarlens}

\subsection*{Your turn}
\begin{itemize}
\raggedright
  \item \emph{Say it:} \emph{klyuch}
  \item \emph{Say it:} \emph{klyuch}, once more
  \item \emph{Say it:} say \emph{klyuch}
  \item \emph{Recall:} say the Russian for a bed, then the Russian for a lamp, then say \emph{klyuch} again
\end{itemize}

\begin{tcolorbox}[breakable,colback=teal!4,colframe=teal!35!black,title={Before you move on}]
What does \ru{ключ} mean? (\textquotedblleft{}A key\textquotedblright{}.) Say it once more.
\end{tcolorbox}
